\documentclass[aspectratio=169, 10pt]{beamer}

% ---------------------------------------------------------------------------
% Call with Trafikverket (noise and vibration), 25 November 2026, 15:00-16:00, in English.
% Core: 4 slides after the title, about 5 minutes. Appendix: for questions.
% Speaker notes, sources and the full question list: trafikverket-meeting-notes.md
% (in the repository, next to this file). Compile with pdfLaTeX and biber.
% ---------------------------------------------------------------------------

\usetheme{metropolis}
\metroset{progressbar=frametitle, numbering=fraction}
\setbeamertemplate{navigation symbols}{}
\usepackage{appendixnumberbeamer}
\usepackage{tikz}
\usetikzlibrary{arrows.meta}
\usepackage[style=authoryear, backend=biber, maxcitenames=2, uniquename=false, uniquelist=false]{biblatex}
\addbibresource{references.bib}

\newcommand{\muted}[1]{{\color{black!55}#1}}
\newcommand{\sv}[1]{{\color{black!55}\textit{(#1)}}}
\AtBeginDocument{\setlength{\leftmargini}{1.3em}}

% a small house at (x,y)
\newcommand{\house}[3][black!45]{%
  \fill[#1] (#2-0.17,#3) rectangle (#2+0.17,#3+0.25);
  \fill[#1] (#2-0.22,#3+0.25) -- (#2,#3+0.43) -- (#2+0.22,#3+0.25) -- cycle;}

\title{Traffic noise, noise protection and long-run human capital in Sweden}
\subtitle{Research plans and data questions}
\author{Petter Lundborg \and Adrian Nieto Castro \and Armando Meier \and Felix Schulz}
\institute{Lund University \quad University of Basel}
\date{Meeting with Trafikverket, 25 November 2026}

\begin{document}

\maketitle

% ===========================================================================
\begin{frame}{What we want to learn}
\begin{columns}[T]
\begin{column}{0.6\textwidth}
\small
\textbf{Question}
\begin{itemize}\setlength{\itemsep}{0pt}
  \item Does road and rail noise in childhood affect learning and later life?
  \item At which ages, and how quickly?
\end{itemize}

\smallskip
\textbf{Why Trafikverket's measures}
\begin{itemize}\setlength{\itemsep}{0pt}
  \item A barrier or a facade measure cuts noise sharply, at a known place and time
  \item This separates noise from everything else near busy roads
\end{itemize}

\smallskip
\textbf{What it gives Trafikverket}
\begin{itemize}\setlength{\itemsep}{0pt}
  \item An estimate of the long-term impact of Trafikverket's noise protection on the children
    who grow up behind it, next to the Swedish evidence on property values and health
    \parencite{lindgrenSoundInvestmentTraffic2021,lindgrenSoundEnvironmentHealth2021}
\end{itemize}
\end{column}
\begin{column}{0.38\textwidth}
\centering
\vspace{1.2em}
\begin{tikzpicture}[scale=1.05]
  % road
  \fill[black!15] (-0.2,0) rectangle (4.6,0.5);
  \draw[white, dashed, thick] (-0.2,0.25) -- (4.6,0.25);
  % noise waves
  \foreach \r in {0.35,0.6,0.85} \draw[red!60!black, thick] (1.2,0.55) ++(30:\r) arc (30:150:\r);
  % wall
  \fill[mDarkTeal] (2.5,0.6) rectangle (4.3,1.05);
  \node[white, font=\scriptsize] at (3.4,0.82) {barrier};
  % houses
  \house[black!45]{0.6}{1.7} \house[black!45]{1.4}{1.9}
  \house[mLightGreen!80!black]{3.0}{1.6} \house[mLightGreen!80!black]{3.8}{1.8}
  \node[font=\scriptsize, align=center] at (1.0,2.75) {loud};
  \node[font=\scriptsize, align=center] at (3.4,2.75) {quieter from\\the year it is built};
  % school
  \fill[black!60] (2.0,-1.2) rectangle (2.9,-0.65);
  \node[white, font=\tiny] at (2.45,-0.93) {school};
  \house[black!45]{0.7}{-1.15} \house[black!45]{3.8}{-1.1}
\end{tikzpicture}
\end{column}
\end{columns}
\end{frame}

% ===========================================================================
\begin{frame}{The idea: two sides of the same road, before and after}
\centering
\begin{tikzpicture}[font=\scriptsize]
  \foreach \x/\lab/\wall in {0/Before the barrier/0, 7.2/After the barrier/1} {
    \begin{scope}[xshift=\x cm]
      \node[font=\small\bfseries] at (2.6,2.65) {\lab};
      \fill[black!15] (0,0) rectangle (5.2,0.5);
      \draw[white, dashed, thick] (0,0.25) -- (5.2,0.25);
      \ifnum\wall=1
        \fill[mDarkTeal] (0.6,0.62) rectangle (4.6,0.82);
        \fill[mLightGreen!30] (0.6,0.9) rectangle (4.6,2.3);
        \node at (2.6,2.1) {protected side: quieter};
      \else
        \fill[black!6] (0.6,0.9) rectangle (4.6,2.3);
        \node at (2.6,2.1) {this side: loud};
      \fi
      \fill[black!6] (0.6,-0.12) rectangle (4.6,-1.5);
      \node at (2.6,-1.3) {other side: still loud};
      \foreach \hx in {1.2,2.2,3.2,4.1} {\house[black!45]{\hx}{1.15} \house[black!45]{\hx}{-0.95}}
    \end{scope}}
  \draw[-{Stealth}, thick, black!60] (5.5,0.4) -- (6.9,0.4);
\end{tikzpicture}

\medskip
\begin{itemize}
  \item Same road, traffic and road works on both sides, but only one side gets quieter
  \item We compare children who grew up on each side, before and after the barrier
\end{itemize}
\end{frame}

% ===========================================================================
\begin{frame}{What we have: the Lastkajen barrier registers}
\begin{columns}[c]
\begin{column}{0.5\textwidth}
\small
\begin{itemize}\setlength{\itemsep}{0pt}
  \item Bullerskydd, Lastkajen products 63 (road, 2,172 records) and 10006 (rail, 1,829)
  \item Linked to every school with coordinates and to 100\,m cells near each barrier
  \item Side checked by hand on aerial imagery
\end{itemize}

\smallskip
\textbf{What is missing}
\begin{itemize}\setlength{\itemsep}{0pt}
  \item No construction year for half of the road and a quarter of the rail records
  \item Road barriers sit on the centreline, with no distance and an unreliable side
  \item Distance from the track for 44\% of rail records
  \item No height for 62\% of road records
\end{itemize}
\end{column}
\begin{column}{0.48\textwidth}
\centering
\textbf{Records by construction year}\\[0.3em]
\begin{tikzpicture}[y=0.0021cm, x=1cm, font=\scriptsize]
  % unknown, split: in the first load (built by then) / registered later
  \fill[mDarkTeal] (-0.42,0) rectangle (-0.02,978);
  \fill[mDarkTeal!40] (-0.42,978) rectangle (-0.02,1114);
  \fill[mLightBrown] (0.02,0) rectangle (0.42,451);
  \fill[mLightBrown!40] (0.02,451) rectangle (0.42,487);
  \node[above, font=\tiny] at (-0.22,1114) {1114};
  \node[above, font=\tiny] at (0.22,487) {487};
  \node[font=\tiny, white, rotate=90] at (-0.22,489) {by 2006};
  \node[font=\tiny, white, rotate=90] at (0.22,225) {by 2019};
  \node[below, align=center, font=\tiny] at (0,0) {unknown};
  % road / rail counts per period: <1996, 1996-2005, 2006-14, 2015-22, 2023+
  \foreach \i/\lab/\road/\rail in {1/{before\\1996}/47/15, 2/{1996--\\2005}/81/349,
                                   3/{2006--\\2014}/174/249, 4/{2015--\\2022}/510/527, 5/{2023--}/246/202} {
    \fill[mDarkTeal] (\i*1.05-0.42,0) rectangle (\i*1.05-0.02,\road);
    \fill[mLightBrown] (\i*1.05+0.02,0) rectangle (\i*1.05+0.42,\rail);
    \node[above, font=\tiny] at (\i*1.05-0.22,\road) {\road};
    \node[above, font=\tiny] at (\i*1.05+0.22,\rail) {\rail};
    \node[below, align=center, font=\tiny] at (\i*1.05,0) {\lab};
  }
  \draw[black!60] (-0.6,0) -- (5.9,0);
  \draw[red!70!black, very thick, rounded corners] (-0.55,-170) rectangle (0.55,1260);
  \fill[mDarkTeal] (2.4,1150) rectangle (2.6,1210); \node[right] at (2.6,1180) {road};
  \fill[mLightBrown] (3.6,1150) rectangle (3.8,1210); \node[right] at (3.8,1180) {rail};
\end{tikzpicture}\\[0.2em]
\muted{\scriptsize Dark: road records loaded in 2006 and rail records entered before April 2019.
These barriers existed by then, but the year they were built is not recorded.}
\end{column}
\end{columns}
\end{frame}

% ===========================================================================
\begin{frame}{What we would like to ask}
\footnotesize
\begin{columns}[T]
\begin{column}{0.36\textwidth}
\textbf{\color{mDarkTeal}1. Noise barriers}\\[0.2em]
\muted{beyond the Lastkajen registers}
\begin{itemize}\setlength{\itemsep}{0pt}
  \item Years for the road barriers loaded in 2006 and the rail barriers entered before 2019
  \item Heights of the older road barriers
  \item Side and distance from the road or track
  \item Does \textit{Byggår} mark completion? Does a rebuild overwrite it?
  \item Older versions of the rail register
  \item Are rail BIS numbers assigned in registration order?
\end{itemize}
\end{column}
\begin{column}{0.28\textwidth}
\textbf{\color{mDarkTeal}2. Facade measures}\\[0.2em]
\muted{the national noise database?}
\begin{itemize}\setlength{\itemsep}{0pt}
  \item Properties considered for mitigation, and whether a measure was carried out
  \item Type of measure, year and cost
  \item Position (the SCB 100\,m grid cell is enough)
\end{itemize}
\end{column}
\begin{column}{0.33\textwidth}
\textbf{\color{mDarkTeal}3. Construction and context}\\[0.2em]
\muted{when, and how loud}
\begin{itemize}\setlength{\itemsep}{0pt}
  \item Start or decision dates for road projects with a project number, and for rail projects
  \item Noise levels calculated before and after a measure
  \item Strategic noise maps from before 2016
  \item Train counts, and road traffic counts before NVDB's history
  \item Barriers planned or proposed but never built
\end{itemize}
\end{column}
\end{columns}

\medskip
\centering
\small Who holds these data, and how could researchers apply for access?
\end{frame}

% ===========================================================================
\appendix
\begin{frame}[standout]
  Appendix
\end{frame}

% ---------------------------------------------------------------------------
\begin{frame}{Design: who counts as exposed}
\begin{columns}[c]
\begin{column}{0.5\textwidth}
\small
\begin{itemize}
  \item Each child is assigned the family's address three years before the barrier, before the
    project could make families move
  \item We follow the child wherever the family moves
  \item We compare cohorts tested before and after the barrier, on the protected side and on the
    opposite side
  \item We expect a step from the first cohort that spends its last school year behind the
    barrier
  \item Longer exposure may add to it
\end{itemize}
\end{column}
\begin{column}{0.48\textwidth}
\centering
\begin{tikzpicture}[x=0.32cm, y=1.6cm, font=\scriptsize]
  \fill[black!6] (-6,-0.2) rectangle (-1.5,1.3);
  \fill[mLightBrown!30] (-1.5,-0.2) rectangle (0.5,1.3);
  \draw[->] (-6.3,0) -- (12,0) node[below left] {cohorts, by test year};
  \draw[->] (-6,-0.2) -- (-6,1.4) node[above] {difference between sides};
  \draw[very thick, red!70!black] (-6,0) -- (1,0) -- (1,0.6) -- (11.5,1.0);
  \node[align=center] at (-3.7,0.85) {tested\\before};
  \node[rotate=90, font=\tiny] at (-0.5,0.85) {works};
  \node[align=left, right] at (1.2,0.45) {step, quieter\\at once};
  \node[align=center, font=\tiny] at (8.5,1.15) {longer exposure};
  \node[below] at (0,-0.2) {built};
\end{tikzpicture}
\end{column}
\end{columns}
\end{frame}

% ---------------------------------------------------------------------------
\begin{frame}{Why construction years matter}
\small
\centering
\begin{tikzpicture}[x=0.175cm, y=0.55cm, font=\scriptsize]
  \draw[->, black!60] (1980,0) -- (2026,0);
  \foreach \y in {1980,1990,2000,2010,2020} {\draw[black!60] (\y,0.1) -- (\y,-0.1) node[below] {\y};}
  \fill[black!25] (1980,0.5) rectangle (2025,0.9);  \node[right] at (2025.3,0.7) {home addresses};
  \fill[black!45] (1996,1.2) rectangle (2022,1.6);  \node[right] at (2025.3,1.4) {grade-9 results};
  \fill[black!45] (2012,1.9) rectangle (2022,2.3);  \node[right] at (2025.3,2.1) {grades 3 and 6};
  \fill[mDarkTeal] (2000,2.8) rectangle (2005,3.2); \node[right] at (2025.3,3.0) {barriers built, birth to grade 9};
  \fill[mDarkTeal!60] (2000,3.5) rectangle (2019,3.9); \node[right] at (2025.3,3.7) {barriers built, four cohorts before and after};
  \fill[mDarkTeal!35] (2014,4.2) rectangle (2021,4.6); \node[right] at (2025.3,4.4) {barriers built, same child before and after};
\end{tikzpicture}

\bigskip
\begin{itemize}
  \item Every comparison is timed from the construction year, so an undated barrier cannot be used
  \item Following children from birth to grade 9 needs barriers built 2000--2005, the years where
    most road records have no year
  \item A first school-level analysis kept 55 schools behind dated barriers and had to drop 58
    behind undated ones
  \item With all years known, open population data suggest about twice as many children behind
    usable rail barriers, and road barriers could enter the birth-to-grade-9 comparison
\end{itemize}
\end{frame}

% ---------------------------------------------------------------------------
\begin{frame}{Why exact positions matter}
\begin{columns}[c]
\begin{column}{0.48\textwidth}
\centering
\begin{tikzpicture}[font=\scriptsize]
  \fill[black!15] (0,0) rectangle (5.5,0.9);
  \draw[red!70!black, very thick, dashed] (0.8,0.45) -- (4.7,0.45);
  \node[red!70!black, below] at (2.75,0.4) {register: on the centreline};
  \fill[mDarkTeal] (0.8,1.25) rectangle (4.7,1.42);
  \node[mDarkTeal, above] at (2.75,1.42) {the barrier itself};
  \draw[{Stealth}-{Stealth}, black!60] (5.0,0.45) -- (5.0,1.33) node[midway, right] {side? distance?};
  \house[mLightGreen!80!black]{1.5}{2.1} \house[mLightGreen!80!black]{3.6}{2.0}
  \house[black!45]{1.5}{-1.0} \house[black!45]{3.6}{-1.1}
  \node[align=center] at (2.6,-1.5) {protected or not? depends on the side};
\end{tikzpicture}
\end{column}
\begin{column}{0.5\textwidth}
\small
\begin{itemize}
  \item Who is protected depends on the side of the road the barrier stands on
  \item Our predicted noise reduction needs the distance from the road and the height
  \item The road register has no distance, and its side field did not hold up in our tests; the
    rail register gives a distance for 44\% of barriers
  \item We placed about 2,600 barriers by hand on aerial imagery, and automatic methods agree
    with that in 95--98\% of cases
  \item Exact positions would replace this and cover the barriers we could not place
\end{itemize}
\end{column}
\end{columns}
\end{frame}

% ---------------------------------------------------------------------------
\begin{frame}{Why facade measures matter}
\begin{itemize}
  \item Insulation is a second treatment. It lowers noise indoors one dwelling at a time and can
    be compared with similar dwellings treated later
    \parencite[as in][]{lindgrenSoundInvestmentTraffic2021}
  \item Insulated homes on the opposite side of a barrier are not really unexposed. Without
    records, our barrier estimates are likely biased toward zero
  \item We link households on the SCB 100\,m grid, so a grid cell per property is enough
  \item Properties that were assessed but not insulated are a natural comparison group
\end{itemize}
\end{frame}

% ---------------------------------------------------------------------------
\begin{frame}[allowframebreaks]{References}
\scriptsize
\printbibliography[heading=none]
\end{frame}

\end{document}
